\section{Projection tests and correlated parameters}
\label{sec:projection}

The only external projection statement used here is the finite
incidence estimate below.  We derive its weighted consequence, a
radial improvement, and the correlated-product estimate needed when
two spatial coordinates have the same rate.  Each test uses the
hereditary finite realizations of Section~\ref{sec:entropy}.

Throughout
this section, $m\to\infty$ denotes the base of the block being tested; a
fixed positive block length is absorbed into this base. A probability law
in a bounded Euclidean set is \emph{Frostman of exponent $d$ at this
scale} if
\[
 \mu(B(v,h))\le m^{o(1)}h^d,\qquad m^{-1}\le h\le 1.
\]
The analogous terminology for directions uses angular intervals. It
suffices to impose these bounds on arbitrarily fine fixed grids of powers
of $m$, allowing an arbitrarily small power loss. Between grid points one
uses monotonicity. A restriction is \emph{dense} if its probability is
$m^{-o(1)}$. At a fixed power scale, \emph{power-small} means bounded by
$m^{-c}$ for some $c>0$; the constant may depend on that scale and on the
strict exponent gaps in the assertion. All limits below keep these gaps
fixed before letting the input losses tend to zero. Bounded enlargements
of bins and boundedly overlapping grids do not affect an exponent.

\subsection{The discretized input and its weighted projection consequence}

We state the external incidence theorem in the form needed here. A
$(\delta,d,C)$-set is a nonempty finite family $P$ of dyadic $\delta$-squares such
that
\[
 |\{p\in P:p\cap B(v,h)\ne\varnothing\}|
 \le C h^d |P|,\qquad \delta\le h\le 1,
\]
with harmless constant enlargements at the mesh scale. Tubes are counted
by their dyadic cells in a bounded slope--intercept chart of line space.
Finitely many charts cover the directions that occur.

\Needspace{5\baselineskip}
\begin{theorem}[Discretized planar Furstenberg estimate]
\label{proj:furstenberg}
Let $0<k\le2$ and $0<\sigma\le1$. For every $\varepsilon>0$ there is
$\eta=\eta(\varepsilon,k,\sigma)>0$ such that the following holds for
sufficiently small dyadic $\delta$. Suppose $P$ is a
$(\delta,k,\delta^{-\eta})$-set of pins in a fixed bounded planar
region, normalized to $[0,1]^2$. For every $p\in P$, let
$\mathcal T(p)$ be a family of $\delta$-tubes meeting $p$, with
$|\mathcal T(p)|\asymp M$, whose directions form a
$(\delta,\sigma,\delta^{-\eta})$-set. Then
\[
 \left|\bigcup_{p\in P}\mathcal T(p)\right|
 \gtrsim_{\varepsilon}
 \delta^{-\min\{k,(k+\sigma)/2,1\}+\varepsilon}M.
\]
Consequently the union has lower exponent at least
$\min\{k+\sigma,(k+3\sigma)/2,1+\sigma\}$ as the nonconcentration
losses tend to zero.
\end{theorem}

This is Theorem~4.1 of Ren--Wang~\cite{RenWang2023}, in the
explicitly cited arXiv version~2, with the notation and nonemptiness
convention of their Definitions~1.3, 2.3--2.4, and 3.1. Their theorem is already a
finite-scale estimate for incident tube pencils. In particular, no
discretization of a Hausdorff-dimension conclusion is being used. The
last assertion follows from $M\gtrsim\delta^{-\sigma+\eta}$, which is
the direction nonconcentration bound at radius $\delta$. Comparable
pencil cardinalities are obtained, when necessary, by a dyadic
pigeonhole. There are only logarithmically many such classes in our
bounded, polynomial-size grids. The dependence of $\eta$ on
$\varepsilon$ is part of the theorem and is retained in every use below.

\begin{lemma}[Weighted projection test]
\label{proj:projection}
Let $\mu$ be a planar Frostman $k$ probability law and let $\zeta$ be a
Frostman $\sigma$ law of directions, where $0<k\le2$ and
$0<\sigma\le1$. Let $E$ be a set of point--direction pairs of dense
$\mu\times\zeta$ measure. For a direction $\omega$, write $E_\omega$
for the corresponding point fiber. If the projections of these fibers
have at most $m^{u+o(1)}$ mesh bins for every retained direction, then
\begin{equation}
 u\ge \min\{k,(k+\sigma)/2,1\}.
 \label{eq:source-15}
\end{equation}
The same conclusion applies to dense point restrictions chosen
separately for each direction. It is enough to have the stated ball
bounds with losses that can be made arbitrarily small.
\end{lemma}

\begin{proof}
First discard direction fibers and pin fibers whose relative masses are
too small, using a threshold with exponent tending to zero sufficiently
slowly. Normalize the restrictions of the reference laws $\mu$ and
$\zeta$ to the retained pin and direction sets, and keep $E$ as a
subset of their product. We continue to denote these reference laws by
$\mu,\zeta$; their Frostman bounds lose only subpower factors, and $E$
still has dense product measure. We now reduce these weighted laws to
the finite sets in Theorem~\ref{proj:furstenberg}.

Sample the reference direction law $\zeta$ \emph{globally}, before
drawing the pencils, so every pencil uses the same sampled population.
For small fixed $0<\rho<\sigma$, take
$n\asymp m^{\sigma-\rho}$ independent samples. Choose a permitted
empirical loss $m^a$ with $a>\rho$, with both exponents below the
input allowance in Theorem~\ref{proj:furstenberg}. The required count
in an interval of radius $h$ is at most $m^a n h^\sigma$.
At $h=m^{-1}$ this threshold is of order $m^{a-\rho}$, whereas
the expected count is at most $n h^\sigma m^{o(1)}$.
Binomial tails therefore give the count bound simultaneously over
the polynomial family of intervals with mesh endpoints. Enlargement
then gives it for all angular intervals. The expected proportion of
samples sharing a mesh bin with another sample is at most
\[
 O\bigl(m^{\sigma-\rho}\sup_\omega
             \zeta(B(\omega,O(m^{-1})))\bigr)
 \le m^{-\rho+o(1)}.
\]
Thus repetitions can be removed at a cost negligible compared with
the dense incidence. The tail failures can also be made negligible
relative to that density. Averaging gives a sample on which dense
incidence survives, with dense direction fibers at a dense set of
pins. Apply the same construction to the reference pin law, using
$m^{k-\rho}$ samples after also requiring $\rho<k$; mesh balls have
a polynomial-size covering family. Restriction to a dense fiber
divides a normalized ball bound by at most a subpower factor.
Finally pigeonhole the pencil cardinalities. Their exponent is
$\sigma-\rho-o(1)$, and the pin exponent is $k-\rho-o(1)$;
the normalized nonconcentration losses remain arbitrarily small.

For each retained point--direction pair, draw the line of constant
projection through that point. The sampled directions number at most
$m^{\sigma+o(1)}$ in total. Hence the projection support hypothesis gives
the global line-count upper bound $m^{\sigma+u+o(1)}$. On the other
hand, Theorem~\ref{proj:furstenberg}, with output error
$\varepsilon$, gives lower exponent
\[
 \sigma-\rho+\min\{k,(k+\sigma)/2,1\}-\varepsilon-o(1).
\]
To make this comparison quantitative, first prescribe an output error,
then choose all sampling, restriction, and Frostman losses below the
input allowance in that theorem, and finally let $m$ tend to infinity.
Sending the output error and $\rho$ to zero proves
\eqref{eq:source-15}.

In a bounded pin chart, the map from a direction to the line through a
given pin is uniformly bi-Lipschitz in its slope coordinate. The passage
between angular bins, lines, and dyadic tubes therefore changes only
bounded covering constants. If angular bounds are initially known only
for balls centered at surviving directions, an arbitrary ball meeting
the support is contained in the ball of twice its radius centered at
one of those directions. This also proves the stated variants.
\end{proof}

The weighted reduction above uses one common sampled direction
population to obtain its line-count upper bound. Its point fibers
may depend on the direction; the dense incidence hypothesis supplies
the needed comparison of laws. In the radial bootstrap below, only
the Furstenberg lower bound is needed, and its direction pencils
will instead be sampled separately at each pin.

\subsection{A radial consequence}

The dispersed/concentrated-tube division below is closely related to
the reciprocal thin-tube bootstrap of Orponen--Shmerkin--Wang
\cite[arXiv:2209.00348v1, Section~1.3, Lemma~2.8, and
Corollary~2.18]{OrponenShmerkinWang2024Radial}.  We prove the required
finite-scale bounds and power-small exceptions here, retaining the
conditional-law hypotheses needed in later applications.

For distinct planar points $x,y$, let $[y-x]$ denote the unoriented
direction of their joining line. For a planar law $\nu$, put
\[
 A_\beta(x,y)
 =\nu\{z\ne x:\operatorname{angle}([z-x],[y-x])\le\beta\}.
\]
Using oriented directions would make only a bounded difference.

\Needspace{3\baselineskip}
\begin{lemma}[Radial nonconcentration away from strips]
\label{proj:radial}
Let $0<s_0\le1$ and let $\nu$ be a planar Frostman $s_0$ law at
resolution $m^{-1}$. Suppose that, for every fixed $\chi>0$, every
strip of width $m^{-\chi}$ has power-small $\nu$ mass, uniformly over
the strip. For every $\sigma<s_0$ and every fixed $h\in(0,1]$, with
$\beta=m^{-h}$, the set of pairs for which
\[
 A_\beta(x,y)>\beta^\sigma
\]
has power-small $\nu\times\nu$ measure. Near-diagonal pairs may be
included in the exceptional set. The assertion holds simultaneously
on any fixed finite grid of scale powers.
\end{lemma}

\begin{proof}
At a fixed power scale $\beta$, all original subpower losses are also
$\beta^{o(1)}$. It is enough to rule out a sequence on which the bad
pair mass is $\beta^{o(1)}$: failure of any power bound would supply
such a subsequence. We use angular grids with bounded overlap, choosing
strict exponent margins so that fixed constants in the definition of
$A_\beta$ are harmless.

We first prove the assertion for $2\sigma<s_0$. At any pin there are
at most $O(\beta^{-\sigma})$ angular bins of mass at least
$\beta^\sigma$. Suppose that pairs in these bins have dense incidence.
Cauchy--Schwarz, applied to target fibers, gives two pins with a dense
common target set. They can be required to be at distance at least
$\beta^\alpha$, for any fixed sufficiently small $\alpha>0$, because
the Frostman estimate makes closer pin pairs power-small. Likewise,
discard the strip of width $\beta^\alpha$ about their joining line.
The strip hypothesis makes the total discarded triple mass
power-small, whereas the common-target triple mass is the square of
the dense incidence and is still subpower.

For a target outside this strip, the two joining directions make an
angle whose sine is at least a constant times $\beta^{2\alpha}$.
Indeed the cross product is the pin separation times the distance to
their joining line, and all distances in the denominator are bounded.
An intersection of two pencil tubes is therefore contained in a ball
of radius $O(\beta^{1-2\alpha})$. The common targets lie in at most
$O(\beta^{-2\sigma})$ such balls. Their total mass is at most
\[
 \beta^{-2\sigma+(1-2\alpha)s_0-o(1)},
\]
which is power-small when $\alpha$ is sufficiently small. This
contradicts dense common incidence and proves the initial range.

For the improvement step, suppose the assertion is already known at
an exponent $\sigma<s_0$, at every fixed power scale. Choose
$\alpha,\eta>0$ so that
\begin{equation}
 3\eta+\alpha<\frac{s_0-\sigma}{2},
 \qquad
 \eta<\frac{\alpha(s_0-\sigma)}2.
 \label{proj:radial-margins}
\end{equation}
We prove it at $\sigma+\eta$. A sufficiently fine finite power grid,
the previous assertion in both orders of a pair, and removal of
power-small exceptions give the following inherited estimates: at
every retained pair the angular masses, at all radii from $\beta$ to
$1$, are bounded by $\beta^{-\tau}$ times the radius to the power
$\sigma$. Here $\tau>0$ is arbitrarily small and is chosen after the
strict margins in \eqref{proj:radial-margins}. Bounds centered at
surviving directions extend to all angular balls by enlargement.

Suppose that angular bins of unrestricted target mass at least
$\beta^{\sigma+\eta}$ carry dense incidence even after this removal.
Call such a bin \emph{dispersed} if its containing tube has
$\nu\times\nu$ mass at least $\beta^{2\sigma+3\eta}$ in pairs
separated by at least $\beta^\alpha$. If the dispersed bins carry
dense incidence, discard pins whose surviving dispersed target fiber
has less than a subpower fraction of the target mass. The normalized
remaining pin law satisfies the planar $s_0$ bounds. At each such pin
$x$, normalize its surviving target fiber and push it forward by
$y\mapsto[y-x]$. The inherited angular bounds make this pin's
direction law Frostman $\sigma$, with arbitrarily small restriction
losses. Sample the pins and then sample each of these direction laws,
using the sampling and collision estimates above; choose comparable
pencil cardinalities. Every selected tube is dispersed, up to a
bounded enlargement, and the direction sets may vary with the pin.
These are the pin set and incident pencils in
Theorem~\ref{proj:furstenberg}. As its input losses tend to zero, it
produces at least
\[
 \beta^{-(s_0+3\sigma)/2+o(1)}
\]
distinct dispersed tubes. Here
$\min\{s_0+\sigma,(s_0+3\sigma)/2,1+\sigma\}
=(s_0+3\sigma)/2$, since $\sigma<s_0\le1$.

A separated subcollection of tubes gives the opposite bound. A pair
of points at separation at least $\beta^\alpha$ belongs to at most
$O(\beta^{-\alpha})$ tubes of this subcollection: the admissible
slopes occupy an interval of length $O(\beta/|z-z'|)$, and for each
slope only boundedly many intercept bins are possible. Integrating
the separated-pair masses of the tubes therefore gives at most
\[
 O(\beta^{-2\sigma-3\eta-\alpha})
\]
tubes. The first strict inequality in
\eqref{proj:radial-margins}, followed by sufficiently small input
losses in Theorem~\ref{proj:furstenberg}, contradicts the lower bound.

It remains to consider nondispersed bins. If such a bin has mass
$b\ge\beta^{\sigma+\eta}$, averaging the separated-pair mass inside
the bin gives a ball of radius $O(\beta^\alpha)$ containing all but
\[
 O(\beta^{2\sigma+3\eta}/b)
 \le O(\beta^{\sigma+2\eta})
\]
of its unrestricted target mass. There are at most
$O(\beta^{-\sigma-\eta})$ heavy bins per pin, so removing these
outside portions costs $O(\beta^\eta)$ in pair mass. Dense incidence
survives.

At a target $y$ with a dense reverse fiber, each surviving pin $x$
now determines a tube carrying unrestricted target mass
$\gtrsim\beta^{\sigma+\eta}$ inside $B(y,C\beta^\alpha)$:
the surviving $y$ lies in the ball chosen for its bin, and doubling
that ball about $y$ retains the same original target mass. If $z$
belongs to the same angular bin from $x$, boundedness of the ambient
region gives
\[
 \left|\det\left(z-y,\frac{y-x}{|y-x|}\right)\right|
 \lesssim\beta.
\]
For fixed $z$, the admissible reverse directions $[x-y]$ therefore
lie within $O(\beta/|z-y|)$ of $[z-y]$, with the trivial bound
when $|z-y|\lesssim\beta$. Apply the inherited reverse angular
estimate to this interval and integrate in $z$. Normalizing the
dense reverse fiber costs only a subpower factor, so the average is
at most
\begin{align*}
 &\beta^{-2\tau}
 \int_{|z-y|\lesssim\beta^\alpha}
       \min\{1,(\beta/|z-y|)^\sigma\}\,d\nu(z)\\
 &\hspace{25mm}\lesssim
       \beta^{\sigma-3\tau+\alpha(s_0-\sigma)}.
\end{align*}
To obtain the last line, split the integral into the ball of radius
$\beta$ and dyadic annuli between $\beta$ and $C\beta^\alpha$,
then use the planar Frostman bound; its exponent exceeds $\sigma$.
Taking $\tau$ sufficiently small, the second inequality in
\eqref{proj:radial-margins} makes this upper bound power-smaller than
$\beta^{\sigma+\eta}$. This is again a contradiction.

The averaging in the nondispersed case uses the original targets
inside each angular bin; thus its discarded mass estimate does not
presuppose that a restricted bin retains its original weight. On a
dense reverse fiber the prior centered angular estimates remain
valid after normalization, with only a subpower loss. These facts
justify both removals used above.

The improvement $\eta$ can be chosen bounded below while $\sigma$
stays a fixed distance below $s_0$. Starting from any exponent below
$s_0/2$ and iterating finitely many times proves the assertion at
every fixed exponent below $s_0$. At each step the desired output
gap is fixed first, then the old angular losses and the
Furstenberg-theorem input losses are chosen. A finite union of the
exceptional sets proves the simultaneous-grid assertion.
\end{proof}

\subsection{Scalar consequences for entropy tangents}

In the order of the coordinates in \eqref{eq:source-9}, write the
full partial derivatives as
\[
 (p,n_1,e,g,t,r)
 =(\partial_C F,\partial_B F,\partial_D F,
   \partial_A F,\partial_Y F,\partial_R F).
\]
Each of these six partials lies in $[0,1]$.
Here $t=\partial_Y F$ and $r=\partial_R F$ are entropy rates; in
\eqref{eq:source-9} the same letters denote truncated parameters.
To avoid imposing a false conclusion when the
input rate vanishes, define, for $z>0$,
\[
 f_z(0)=0,\qquad f_z(l)=\min\{1,z+l\}\quad(l>0).
\]

\Needspace{3\baselineskip}
\begin{proposition}[Scalar projection inequalities]
\label{proj:scalar}
At almost every point of the full alignment there is a tangent with
constant full partials, representing the derivatives along the
alignment, for which
\begin{equation}
 e\ge f_t(p),\qquad g\ge f_t(n_1),\qquad
 n_1\ge f_r(p),\qquad g\ge f_r(e),
 \label{eq:source-16}
\end{equation}
whenever the corresponding subscript is positive.
\end{proposition}

\begin{proof}
On the full alignment the six coordinate forms are
\[
 C,\quad C+R,\quad C+Y,\quad C+Y+R,\quad Y,\quad R.
\]
They are pairwise nonproportional and positive on a common vector.
The peeling rule of Lemma~\ref{ent:peeling}, applied successively to
extreme rays of the cone generated by these forms, makes each
coordinate a singleton summand in a tangent. Hence the tangent is
affine with the stated full partials.

We give the normal-velocity/time test in detail. Choose a block length
$\ell>0$ and a parent state $z_0$ on $A=B+Y$, with each of the
other three slacks in \eqref{eq:source-10} larger than $\ell$.
Keep $C,D,R$ fixed. Writing $\mathbf e_i$ for the coordinate vectors,
compare these four refinements of $z_0$:
\[
\begin{array}{c|c}
 \text{state}&\text{entropy increment from }z_0\\ \hline
 z_{\rm sp}=z_0+\ell(\mathbf e_B+\mathbf e_A)&\ell(n_1+g)\\
 z_{\rm par}=z_0+\ell\mathbf e_Y&\ell t\\
 z_{\rm joint}=z_0+\ell(\mathbf e_B+\mathbf e_A+\mathbf e_Y)
      &\ell(n_1+g+t)\\
 z_{\rm out}=z_0+\ell(\mathbf e_A+\mathbf e_Y)&\ell(g+t).
\end{array}
\]
The spatial and output refinements preserve $A=B+Y$; the parameter
and joint states lie on its admissible side. The chosen slacks make
all four states admissible. Their increments are those of the affine
tangent.

Inside a common parent cell, subtract the known bin centers and
normalize by the parent units. Let $v,w$ be the two skew-normal
spatial entries in the starting parameter frame, and let $a$ be the
normalized time increment. The matched equality makes the position
unit the product of the velocity and time units, so the output
measurement is the scalar shear $w+av$. The spatial state measures
$(v,w)$ and the parameter state measures $a$.
Lemma~\ref{ent:finite-states} identifies their joint labels with
$z_{\rm joint}$ up to bounded ambiguity: the fine time label permits
rebasing with the finely measured velocity. Center-dependent terms
are known translations, and \eqref{eq:source-10} bounds the
remaining errors.

Lemmas~\ref{ent:regularization} and~\ref{ent:tangent-realization}
turn these increments into hereditary support and weight bounds.
At every fixed block subresolution, the spatial law is planar
Frostman $n_1+g$, the coefficient law is Frostman $t$, and their
joint exponent is $n_1+g+t$.

We make the comparison with the actual marginal laws explicit.
Let $\mu_1,\mu_2$ be the marginals of the same joint law, with
support exponents $\alpha,\beta$ and maximal joint-bin weight
$m^{-\alpha-\beta+o(1)}$. In any incidence of joint mass
$\rho=m^{-o(1)}$, discard marginal bins with reference weights
below $m^{-\alpha-o(1)}$ and $m^{-\beta-o(1)}$, choosing the
threshold slack to tend to zero sufficiently slowly. The support
bounds make their marginal masses, and hence the joint mass incident
to them, $o(\rho)$. Since the paired labels encode the joint cells
up to bounded multiplicity, at least
$\rho m^{\alpha+\beta-o(1)}$ pairs remain. Each has
$\mu_1\times\mu_2$ weight at least
$m^{-\alpha-\beta-o(1)}$. Their total product mass is therefore
at least $\rho m^{-o(1)}=m^{-o(1)}$. For the present labels take
$\alpha=n_1+g$ and $\beta=t$. This proves dense incidence for
the actual Frostman marginals, also after the required subpower
joint restrictions.

Conditioning $z_{\rm out}$ on $z_{\rm par}$ leaves exponent $g$,
so the projections occupy at most $m^{g+o(1)}$ bins for typical
fine time labels, where $m$ absorbs the block length. The parent
and time-bin trims can be made negligible relative to this dense
incidence. If $n_1,t>0$,
Lemma~\ref{proj:projection} gives
\[
 g\ge\min\{n_1+g,(n_1+g+t)/2,1\}.
\]
If $g<1$, the first and third entries of this minimum exceed $g$,
so $g\ge n_1+t$. If $g=1$, the desired inequality is automatic.
This proves $g\ge f_t(n_1)$; the case $n_1=0$ is also automatic
by the definition of $f_t(0)$.

The other three tests use the following triples and shear maps:
\[
\begin{array}{c|c|c}
 \text{cutoffs}&\text{matched face}&\text{projection}\\ \hline
 (C,D,Y)&D=C+Y&z+ax\\
 (C,B,R)&B=C+R&y-ax\\
 (D,A,R)&A=D+R&w-az.
\end{array}
\]
Here the unused cutoffs are fixed with strict slack, and $a$ denotes
the parameter increment in the corresponding normalized units. The
same joint-count and conditional-support calculation proves the
remaining assertions. The shear update
\eqref{ent:shear-update} controls the labels at the unused cutoffs; it
does not require the normalized spatial law to be independent of
the coefficient law. Only dense product incidence is used.

The tangent is affine throughout its admissible domain. Placing a
test on its matched face with slack at the other constraints leaves
all six rates unchanged. The zero-input and saturated-output cases
are automatic; in every other case the projection minimum gives the
linear inequality that the output rate is at least the sum of the
input and parameter rates. Thus each displaced test proves the
assertion for the original constant rates.
\end{proof}

\begin{proposition}[Conditional parameter profile at an increasing tie]
\label{proj:parameter-profile}
Suppose $\lambda>0$. At a generic point of the calibrated line one
may take a further joint tangent in which the conditional parameter
profile, referenced to zero at the center, is
\begin{equation}
 K(Y,R)=sY+\theta(R-\lambda Y)_+,
 \qquad 0\le\theta\le s/\lambda.
 \label{eq:source-17}
\end{equation}
Consequently $t\ge s$ below the tie, while $r\ge\theta$ and
$t\ge s-\lambda\theta$ above it. The simultaneous parameter
advance $(Y,R)=(u,\lambda u)$ has rate at least $s$.

Let $\Gamma$ be the plane $R=\lambda Y$ inside full alignment. If
$\lambda\ne1$, the spatial partials can simultaneously be made
singleton constants at generic points of $\Gamma$. There
\eqref{eq:source-16} holds with time subscript $s$, and also with
tilt subscript $\theta$ when $\theta>0$.
\end{proposition}

\begin{proof}
The predictability statement of Section~\ref{sec:entropy} gives
$K(Y,R)=sY$ for $R\le\lambda Y$. Apply the increasing-coordinate
trace rule proved with Lemma~\ref{ent:signed-envelope} to the two
conditional parameter coordinates. A generic tangent along their
increasing tie is affine on each of the two ordering half-planes.
Continuity along $R=\lambda Y$ forces its upper-side expression to
be $sY+\theta(R-\lambda Y)$. Monotonicity in $R$ and $Y$ gives
$\theta\ge0$ and $s-\lambda\theta\ge0$, respectively. This proves
\eqref{eq:source-17}. Conditional parameter increments lower-bound
the corresponding full-state increments, giving all the asserted
rate bounds without an independence assumption.

On $\Gamma$ the spatial forms are
\[
 C,\qquad C+\lambda Y,\qquad C+Y,
 \qquad C+(1+\lambda)Y.
\]
For $\lambda>0$, $\lambda\ne1$, they are distinct,
nonproportional forms, and the only remaining proportional group
consists of the two parameters. Peeling separates the four spatial
singletons, whose rates are constant throughout this tangent.
At a generic full-alignment point on the lower side $R<\lambda Y$,
take a further realized tangent. It has the same four spatial rates,
and the conditional-parameter increment bound gives its time rate
at least $s$. Proposition~\ref{proj:scalar} therefore gives the time
inequalities with subscript $s$. At a generic full-alignment point
on the upper side $R>\lambda Y$, the same construction gives tilt
rate at least $\theta$ and the tilt inequalities with that subscript.
The two further tangents may have different parameter partials, but
both conclusions concern the original four constant spatial rates.
They therefore hold simultaneously at the tie.
\end{proof}

\subsection{Products of correlated parameter differences}

At $\lambda=1$ the forms $B$ and $D$ coincide on $\Gamma$.
Accordingly the preceding scalar tests alone do not separate their
rates. We first prove the parameter fact needed for this case.

\begin{lemma}[Product coefficients]
\label{proj:product}
Let $\nu$ be a law of pairs $d=(r_d,t_d)$ in a bounded planar set.
Assume its time marginal is Frostman $s$, its tilt marginal is
Frostman $\theta$ for some $\theta>0$, and $0<s\le1$. Assume also
that its simultaneous two-coordinate refinements have an affine
hereditary profile of dimension at least $s$: in typical
intermediate cells the normalized unrestricted increments satisfy
the Frostman $s$ bounds at all subsequent fixed block
subresolutions. For a pin $f=(r_f,t_f)$ define
\[
 Q_f(d)=(r_d-r_f)(t_d-t_f).
\]
The following alternatives suffice at every fixed finite grid of
resolutions.
\begin{enumerate}[label=(\roman*)]
\item If some strip of width $m^{-\chi}$, with fixed $\chi>0$,
has dense $\nu$ mass, restrict $\nu$ to that strip. At any fixed
resolution $m^{-h}$ with $0<h<\chi$ sufficiently small, the
coefficients $Q_f(d)$ have Frostman exponent $s$, up to subpower
loss, on dense pairs away from the time diagonal.
\item If every strip of each fixed power width has power-small
mass, then for every $\sigma<s$ and every fixed $h\in(0,1]$,
pairs lying in coefficient bins of conditional mass greater than
$m^{-h\sigma}$ have power-small total mass. After removal of
these exceptions on a sufficiently fine finite grid, the pinned
coefficient laws satisfy Frostman bounds with exponent
arbitrarily close to $s$ and arbitrarily small power loss.
\end{enumerate}
\end{lemma}

\begin{proof}
\textbf{The dense-strip alternative.}
Suppose first that a strip of width $m^{-\chi}$ has dense mass.
Its slope $a=dr/dt$ satisfies
\[
 m^{-o(1)}\le |a|\le m^{o(1)}.
\]
Indeed, if $|a|\le m^{-\varepsilon}$ along a subsequence, its
intersection with the bounded parameter region lies in a tilt
interval of length $O(m^{-\varepsilon}+m^{-\chi})$. The positive
tilt Frostman exponent makes its mass power-small. If
$|a|\ge m^\varepsilon$, including the vertical case, use the time
marginal instead. Since this applies to every fixed
$\varepsilon>0$, it proves the claim. Conditioning on the strip
costs only a subpower factor in either marginal bound.

On the strip write
\[
 r=at+b+e,
 \qquad |e|\le m^{-\chi}\sqrt{1+a^2}=m^{-\chi+o(1)}.
\]
Uniformly in its pins and targets,
\[
 Q_f(d)=a(t_d-t_f)^2+O(m^{-\chi+o(1)}).
\]
The time Frostman bound permits removal of pairs with
$|t_d-t_f|<m^{-o(1)}$, with the cutoff decreasing sufficiently
slowly that the lost mass is negligible compared with every dense
incidence being retained. On each sign branch of $t_d-t_f$ the
quadratic has derivative of magnitude between $m^{-o(1)}$ and
$m^{o(1)}$. At radii $\rho\ge m^{-h}$, with fixed $h<\chi$,
the inverse image of a coefficient interval is therefore contained
in at most two time intervals of length $m^{o(1)}\rho$. The time
marginal bound proves (i), also at all the intermediate radii
needed for a projection test.

\paragraph{Power-light strips and the conditional query law.}
For (ii), the planar law is Frostman $s$, by the simultaneous
refinement hypothesis. Lemma~\ref{proj:radial} therefore applies.
Fix $\sigma<s$ and $h>0$. Suppose that pairs landing in bins of
$Q_f$ of conditional mass at least $m^{-h\sigma}$ carry dense
incidence $E$. Condition the product law on this incidence. For
each pin the conditioned output occupies at most
$O(m^{h\sigma})$ bins, so, with $P_h$ the output bin,
\begin{equation}
 \Ent(P_h\mid f)\le h\sigma\log m+O(1).
 \label{proj:product-upper}
\end{equation}
All entropies in the rest of this proof are under the conditioned
law unless explicitly described as unrestricted.

Choose $\sigma<s'<s$ and divide $[0,h]$ into $J$ equal blocks
of length $\eta=h/J$. The integer $J$ will be large but fixed.
Let $D_j$ be the target parameter cell at depth $j\eta$, and
let $P_j$ be the nested bin of $Q_f$ at that depth. For
$j\ge3$ we claim that, apart from a negligible conditioned mass,
\begin{equation}
 \Ent(P_{j+1}\mid f,D_j)
 \ge (s'-\varepsilon)\eta\log m-o(\log m),
 \label{proj:product-block}
\end{equation}
where $\varepsilon>0$ is arbitrarily small. We verify carefully
the conditioning and direction estimates behind this assertion.

Use the hereditary parent estimates on this fixed block grid under
the original target law $\nu$, before conditioning on $E$. Choose
their errors and parent-trim thresholds as in
Lemma~\ref{ent:regularization}, so that the omitted original cell
mass is negligible compared with the dense incidence being retained.
For a remaining value $C$ of $D_j$, let $c$ be its center and let
$\mu_C$ be the normalized original restriction $\nu|_C$, expressed
in the cell's increment units. This \emph{unrestricted} law is
Frostman $s$ at all the required block subresolutions. The following
density trims compare the conditioned query and target fibers to
this law. The map
$\psi_C(v)=c+m^{-j\eta}v$ returns an increment to its physical
target cell. Define the pulled-back incidence and its target fibers by
\[
 E_C=\{(f,v):(f,\psi_C(v))\in E\},\qquad
 E_{C,f}=\{v:(f,v)\in E_C\}.
\]
The pin law $\nu$ remains in the original parameter coordinates.
Set
\[
 \alpha_C=(\nu\times\mu_C)(E_C).
\]
After discarding cells with too small an incidence density,
$\alpha_C=m^{-o(1)}$ in cells carrying almost all of the
conditioned mass. To see this, the total incidence in cells with
$\alpha_C<\rho$ is at most $\rho$; choose the subpower threshold
$\rho$ much smaller than the original dense incidence. The query
law in such a cell is
\begin{equation}
 q_C(df)=\frac{\mu_C(E_{C,f})}{\alpha_C}\,\nu(df)
 \le m^{o(1)}\nu(df).
 \label{proj:query-domination}
\end{equation}
Another density trim ensures $\mu_C(E_{C,f})=m^{-o(1)}$ for almost
all contributing queries. The same reasoning applies after any
of the finitely many required power-small exceptional pair sets
has been removed.

\paragraph{Radial control of the query directions.}
The gradient of $Q_f$ at $c$ is
\[
 \nabla Q_f(c)=(t_c-t_f,r_c-r_f).
\]
Thus its direction is the radial direction from $f$ to $c$ with
the coordinates interchanged. Remove the exceptional pairs from
Lemma~\ref{proj:radial}, in both orders and at a fine grid of
angular powers, using exponent $s'$, for the physical pairs
$(f,\psi_C(v))$. This removes a negligible
fraction of the dense incidence. Direction estimates may be
transferred from an actual surviving target $d\in C$ to the
center $c$. For comparison pins at distance at least
$m^{-\eta}$ from $c$, this changes the direction by
\[
 O(m^{-j\eta}/m^{-\eta})
   =O(m^{-(j-1)\eta}),
\]
which is below the block angular mesh for $j\ge3$. Comparison
pins closer than $O(m^{-\eta})$ have mass at most
$m^{-\eta s+o(1)}$ by planar Frostman. The surviving queried
pin itself can be required to be at distance at least
$m^{-o(1)}$ from the target, so its direction also changes by
less than the block mesh. For a ball centered at a surviving
query direction, choose its corresponding good target in $C$
and use its radial estimate. Enlarging the angular ball by the
above errors bounds all the far comparison pins; the near ones
have just been bounded separately. Re-centering balls extends
the conclusion to every angular interval. In view of
\eqref{proj:query-domination}, the law of gradient directions
under $q_C$ is consequently Frostman $s'$, with any prescribed
small power loss down to the block mesh. The grids and prior
losses are chosen sufficiently fine and small for that purpose.

\paragraph{Collision bounds inside a target cell.}
Here is the elementary projection-energy calculation in the
normalized cell. Let $v,v'$ be independent with law $\mu_C$,
and let $\omega$ have the gradient-direction law. A projection
collision at mesh $m^{-\eta}$ requires $\omega$ to lie in
boundedly many angular intervals of total length
$O(m^{-\eta}/|v-v'|)$. Therefore its averaged probability is
bounded by
\begin{align}
 &m^{\varepsilon_0}
 \iint\min\{1,(m^{-\eta}/|v-v'|)^{s'}\}
                       \,d\mu_C(v)d\mu_C(v')
 \notag\\
 &\hspace{35mm}\le m^{-\eta s'+\varepsilon_1}.
 \label{proj:collision}
\end{align}
The last estimate follows by summing annuli and using the
Frostman $s$ bound, with $s>s'$. Both
$\varepsilon_0,\varepsilon_1$ can be chosen arbitrarily small
compared with $\eta$. Markov's inequality gives this collision
bound, with another arbitrarily small loss, for almost all
queries. Restriction to $E_{C,f}$ multiplies the collision
probability by at most $\mu_C(E_{C,f})^{-2}=m^{o(1)}$. Thus the
estimate applies to the target fibers actually present in the
conditioned incidence, rather than just to unrestricted
increments.

The gradient magnitude is at least $m^{-o(1)}$ after the
near-diagonal removal. Normalizing it changes the mesh to
$m^{-\eta+o(1)}$, an admissible subpower change in
\eqref{proj:collision}. Taylor expansion gives, for bounded $v$,
\[
 Q_f(c+m^{-j\eta}v)
 =Q_f(c)+m^{-j\eta}\nabla Q_f(c)\cdot v
      +m^{-2j\eta}v_rv_t.
\]
The remainder is below the physical next resolution
$m^{-(j+1)\eta}$. Hence a bin of the nonlinear output is
covered by boundedly many bins of the linearized output at the
relevant normalized mesh. Finally, for bin probabilities $a_i$,
\[
 -\sum_i a_i\log a_i\ge-\log\sum_i a_i^2.
\]
The collision estimate yields \eqref{proj:product-block}.
The negligible removed conditioned mass is restored using
entropy concavity. All bounds are averaged with respect to the
actual conditional query laws, so no independence after
conditioning has been assumed.

\paragraph{Summing the output entropy increments.}
Because $Q_f$ is uniformly Lipschitz on the bounded parameter
region, $D_j$ together with $f$ determines $P_j$ up to bounded
ambiguity. Nested output bins also give
\begin{align}
 \Ent(P_{j+1}\mid f)-\Ent(P_j\mid f)
 &\ge
 \Ent(P_{j+1}\mid f,D_j)-\Ent(P_j\mid f,D_j).
 \label{proj:product-chain}
\end{align}
Indeed the left side is
$\Ent(P_{j+1}\mid f,P_j)$, and conditioning further on $D_j$
can only decrease it. The subtracted right-hand term is $O(1)$.
Sum \eqref{proj:product-chain} for $j=3,\ldots,J-1$ and use
\eqref{proj:product-block}. This gives
\[
 \Ent(P_h\mid f)
 \ge (J-3)(s'-\varepsilon)\eta\log m-o(\log m).
\]
Choose $J$ so large, and then $\varepsilon$ and all other losses
so small, that $(1-3/J)(s'-\varepsilon)>\sigma$. This contradicts
\eqref{proj:product-upper}.

We have ruled out a subpower sequence of bad pair masses, and
therefore proved a power-small bound at each fixed coefficient
scale. Remove these exceptions on a finite fine grid. The
conditional mass bound at the bin containing any surviving
coefficient gives the centered Frostman bound there; an interval
meeting the retained support is covered by a bounded enlargement
centered at one of its points. Normalizing dense pin fibers and
interpolating between grid powers costs an arbitrarily small
power. This proves (ii).
\end{proof}

\subsection{The tied spatial pair}

We now return to the coincident spatial forms.  The middle pair
$(B,D)$ remains one joint measurement.  Overlaps of full spatial
fibers will produce scalar projections with coefficients supplied
by Lemma~\ref{proj:product}; no separate rates for the two tied
coordinates are assumed.

\begin{proposition}[The equal-slope tie]
\label{proj:tied}
Suppose $\lambda=1$ and $\theta>0$ in
\eqref{eq:source-17}. At a generic point of $\Gamma$ there is a
tangent with singleton spatial rates $p,g$, an affine diagonal
rate $H$ for the pair $(B,D)$, and an affine diagonal parameter
rate. These rates satisfy
\begin{equation}
 g\ge f_s(p),\qquad H\ge f_s(p).
 \label{eq:source-18}
\end{equation}
\end{proposition}

\begin{proof}
On $\Gamma$ the coordinate forms are $C$, $C+Y$ twice,
$C+2Y$, and $Y$ twice. Peeling separates the singleton groups
$C,A$ and the paired groups $(B,D)$ and $(Y,R)$; each group
has an affine restriction to its diagonal. This gives the
claimed form. Write $J(B,D)$ for the separated middle-pair summand
of this actual tangent. On lower-side full alignment, $R<Y$ is
equivalent to $B<D$. Further full-alignment tangents there preserve
$p$ and have time rate at least $s$, so the scalar time test gives
$\partial_DJ\ge f_s(p)$ almost everywhere in $B<D$. For $h>0$,
monotonicity in $B$ and integration of this $D$-bound give
\[
 Hh=J(b+h,b+h)-J(b,b)
 \ge J(b,b+h)-J(b,b)\ge h f_s(p).
\]
The integration first holds on almost every vertical segment;
continuity includes every such segment and its boundary endpoint.
This proves the inequality for $H$ without assigning separate
$B,D$ partials at the tie. If $p=0$,
both assertions in \eqref{eq:source-18} are immediate. Assume
from now on that $p>0$.

Take a block in a typical fully aligned parent on the tie. In normalized units
write the spatial entries in its starting frame as $x,y,z,w$,
and fine parameter pairs as $d=(r_d,t_d)$ with law $\nu$.
The parameter diagonal rate is at least $s$. Heredity makes
its law Frostman of that diagonal dimension, also inside
intermediate parameter cells at every later fixed
subresolution. The time marginal is Frostman $s$, and the
tilt marginal has exponent at least $\theta>0$, by
\eqref{eq:source-17} and the conditional-parameter lower
bound. Thus the parameter laws satisfy the hypotheses of
Lemma~\ref{proj:product}.

Denote the parent by $z_0$ and the block length by $\ell>0$.
With $\mathbf e_i$ denoting the coordinate vectors, set
\[
\begin{aligned}
 z_{\rm sp}&=z_0+\ell(\mathbf e_C+\mathbf e_B+\mathbf e_D+\mathbf e_A),\\
 z_{\rm par}&=z_0+\ell(\mathbf e_Y+\mathbf e_R),\\
 z_{\rm joint}&=z_{\rm par}
                  +\ell(\mathbf e_C+\mathbf e_B+\mathbf e_D+\mathbf e_A).
\end{aligned}
\]
The spatial increment of $z_{\rm sp}$ has exponent $p+H+g$.
Group additivity gives the same exponent for $z_{\rm joint}$
conditional on $z_{\rm par}$. The marginal-weight comparison in
the proof of Proposition~\ref{proj:scalar}, with spatial exponent
$p+H+g$ and the parameter diagonal exponent, therefore gives dense
incidence under the product of the actual spatial and parameter
marginals. Given the fine
parameter label $d$, write the spatial entries in its frame as
$x,y_d,z_d,w_d$. The two further states
\[
 z_{\rm par}+\ell(\mathbf e_C+\mathbf e_A),\qquad
 z_{\rm par}+\ell(\mathbf e_B+\mathbf e_D+\mathbf e_A)
\]
have increments $\ell(p+g)$ and $\ell(H+g)$ from $z_{\rm par}$.
The remaining refinement from the first of these states to
$z_{\rm joint}$ costs $\ell H$. All these states satisfy
\eqref{eq:source-10}; in particular the fine parameters provide
the slack for these partial spatial refinements.

Lemma~\ref{ent:finite-states} rebases the full spatial labels using
the known fine parameters, with bounded ambiguity.
Lemmas~\ref{ent:regularization} and~\ref{ent:tangent-realization},
applied to these state pairs at every required block subresolution,
make the conditional $(x,w_d)$ law planar Frostman $p+g$.
Its additional $(y_d,z_d)$ fibers have support upper exponent $H$,
and the output triples $(y_d,z_d,w_d)$ have support upper exponent
$H+g$. We use these conditional laws before any overlap restriction.

We spell out how the dense spatial overlaps carry the weights
needed for Lemma~\ref{proj:projection}. This also ensures that
subsequent choices of parameters do not replace a weighted
projection problem by a mere cardinality assertion. Let $m'$
be the base of the positive-length block presently used; it
may be a fixed fractional block. In a typical parent, let
$S_f$ be the coarse-frame spatial bins incident to a fine
parameter bin $f$. Coarse-frame and fine-frame full spatial
bins determine one another up to bounded coverings when $f$
is given. The exact truncated parameters are used in
subtracting center-dependent translations, so large absolute
centers create no new error.

Put $L=p+H+g$. The common spatial universe is the support of the
spatial label at $z_{\rm sp}$, of size at most $(m')^{L+o(1)}$.
For almost every
tested $f$, $S_f$ has size $(m')^{L+o(1)}$, and every retained
joint cell has pre-trim conditional weight at least
$(m')^{-L-o(1)}$ given $f$. These conclusions follow from the support
and maximal-weight bounds for $z_{\rm joint}$ conditional on
$z_{\rm par}$, followed by discarding cells below the
reciprocal-support threshold. The lower weights refer to this
conditional law before that final discard. The exceptional fraction
can be made negligible compared with any prescribed dense
parameter restriction, by taking the threshold slack to zero
sufficiently slowly. In particular the entire sets $S_f$
remain available after such a restriction.

For any dense restriction of the parameter law, Cauchy--Schwarz
in the common spatial universe gives
\begin{align*}
 \iint |S_f\cap S_d|\,d\nu(f)d\nu(d)
 &=\sum_x\left(\int\mathbf1_{S_f}(x)\,d\nu(f)\right)^2\\
 &\ge
 \frac{\left(\int |S_f|\,d\nu(f)\right)^2}
      {|\text{spatial universe}|}
 \ge (m')^{L-o(1)}.
\end{align*}
Here $x$ in the sum denotes a full spatial cell, not only its
first coordinate. Since the overlap upper bound is
$(m')^{L+o(1)}$, a dense set of parameter pairs has overlaps
of full exponent $L$. Power-small exceptional parameter pairs
can be removed before selecting these overlaps.

For one such pair $f,d$, put
$\Delta r=r_d-r_f$ and $\Delta t=t_d-t_f$. The shear update
gives, to mesh accuracy,
\begin{align}
 y_d&=y_f-\Delta r\,x,
 &z_d&=z_f+\Delta t\,x,
 \notag\\
 w_d-\Delta t\,y_d+\Delta r\,z_d
     &=w_f+\Delta t\Delta r\,x.
 \label{proj:overlap-identity}
\end{align}
The last identity follows by substituting the first two into
the full shear update, including its mixed term. Bounded bin
errors produce only bounded mesh errors in normalized units.

There are at most $(m')^{p+g+o(1)}$ $(x,w_f)$ bins in the
overlap. Discard those carrying fewer than
$(m')^{H-o(1)}$ additional full cells, with the slack chosen
sufficiently slowly. They account for a negligible fraction
of an overlap of exponent $L$. Call the remaining input bins
rich. Inside one rich input bin, the different full cells map
to different triples $(y_d,z_d,w_d)$ up to bounded
multiplicity: once $x$ is fixed to mesh accuracy, the frame
change can be inverted from that triple. Thus a scalar value
on the right of \eqref{proj:overlap-identity}, attained in
one rich input bin, uses at least $(m')^{H-o(1)}$ output
triples. Scalar values separated by a sufficiently large
mesh constant use disjoint triples, because the left side
of that identity is determined to mesh accuracy by a triple.
The total output support has size at most
$(m')^{H+g+o(1)}$. Consequently the scalar projection
\[
 (x,w_f)\longmapsto w_f+\Delta r\Delta t\,x
\]
of the rich subset has at most $(m')^{g+o(1)}$ bins.

These rich subsets are dense for the original conditional
$(x,w_f)$ law. Indeed they retain $(m')^{L-o(1)}$ joint cells,
each with pre-trim conditional weight at least
$(m')^{-L-o(1)}$. Their total weight is therefore
$(m')^{-o(1)}$, and their image consists of the corresponding
coarse-accuracy input bins. The planar Frostman bound is the
one before this final restriction, obtained from the matched
$C,A$ advances. The rich restriction costs only a subpower
factor. This is the required passage from overlap counts to
weighted point fibers.

We now obtain enough coefficient directions. If $\nu$ has
dense mass in a strip of width $m^{-\chi}$ for some fixed
$\chi>0$, first restrict the parameters to this strip and
repeat the overlap argument with the full spatial sets of
the typical parameter bins. Use a fixed positive fractional
block ending at depth $h<\chi$ sufficiently small.
Lemma~\ref{proj:product}(i) gives a Frostman $s$ law of
$\Delta r\Delta t$ on dense pairs. The removed time-diagonal
pairs have negligible mass compared with the dense overlaps.
All spatial rate, fiber, and weight statements above apply
at this fractional depth with $m'=m^h$.

If there is no such strip restriction along the realizing
sequence, pass to a subsequence on which every strip of each
fixed power width has power-small mass. This dichotomy is
obtained by inspecting the maximum strip mass on countably
many fixed power grids; bounded mesh enlargements suffice.
Lemma~\ref{proj:product}(ii) gives coefficient exponents
arbitrarily close to $s$ off power-small exceptional pairs.
Remove these pairs on the finite fine grid of coefficient
resolutions before choosing a pin with dense overlap.

In either case, averaging selects $f$ with densely many
targets $d$ that have the coefficient Frostman bounds and
the rich projected-input subsets just constructed. Parameter
binning changes the coefficient only by mesh order. If
several target bins produce the same coefficient bin, choose
one corresponding rich input subset for that bin; each has
the required dense mass for the same pre-trim conditional
$(x,w_f)$ law. Thus the coefficient marginal together with
these rich subsets supplies dense product incidence for
Lemma~\ref{proj:projection}. The map from a bounded scalar
coefficient to the associated projection direction is
bi-Lipschitz in a bounded chart.

Applying that lemma and letting all coefficient and
restriction losses tend to zero yields
\[
 g\ge\min\{p+g,(p+g+s)/2,1\}.
\]
Since $p>0$, either $g=1$, or the first and third terms
exceed $g$ and the middle term gives $g\ge p+s$. In both
cases $g\ge f_s(p)$, completing the proof.
\end{proof}
